% Problem record op_0c86d9293aba3b01. This TeX file is derived from
% database/problems_json/op_0c86d9293aba3b01.json by scripts/sync-tex.mjs; edit the JSON record
% and rerun the script rather than editing this file.
\section{Finite nontrivial LU moduli of AME states}
\label{sec:op_0c86d9293aba3b01}

\paragraph{Problem.}
Do there exist integers $N,d\geq 2$ for which the absolutely maximally
entangled states of $N$ qudits of local dimension $d$ form finitely many, but
more than one, local-unitary equivalence classes?  A normalized vector
$|\psi\rangle\in(\mathbb{C}^{d})^{\otimes N}$ is an $\operatorname{AME}(N,d)$
state when every subsystem of at most half the parties is maximally mixed:
\begin{equation}
  \operatorname{Tr}_{S^{c}}\!\left(|\psi\rangle\!\langle\psi|\right)
  =\frac{I_{d^{|S|}}}{d^{|S|}}
  \quad\text{for every }S\subseteq\{1,\ldots,N\}
  \text{ with }|S|\leq\left\lfloor\frac{N}{2}\right\rfloor .
  \label{eq:p43-ame-condition}
\end{equation}
For normalized states satisfying Eq.~\eqref{eq:p43-ame-condition}, define
local-unitary equivalence by
\begin{equation}
  |\psi\rangle\sim_{\mathrm{LU}}|\phi\rangle
  \quad\Longleftrightarrow\quad
  |\psi\rangle=(U_1\otimes\cdots\otimes U_N)|\phi\rangle
  \quad\text{for some }U_1,\ldots,U_N\in U(d).
  \label{eq:p43-lu-equivalence}
\end{equation}
If $\mathcal{A}_{N,d}$ denotes the set specified by
Eq.~\eqref{eq:p43-ame-condition}, determine whether the quotient under
Eq.~\eqref{eq:p43-lu-equivalence} can satisfy
\begin{equation}
  \mathfrak{M}_{N,d}:=\mathcal{A}_{N,d}/\!\sim_{\mathrm{LU}},
  \qquad
  1<\lvert\mathfrak{M}_{N,d}\rvert<\infty .
  \label{eq:p43-finite-nontrivial-moduli}
\end{equation}

\subsection*{Status}

Unsolved

\subsection*{Source}

Rajchel-Mieldzio\'c et al. explicitly ask whether an AME family can have a
finite number of local-unitary equivalence classes greater than one
\sourcecite{ref:p43-review}{RBR+26}.

\subsection*{Progress}

\begin{itemize}
  \item For four parties, all $\operatorname{AME}(4,3)$ states lie in one LU
  class, whereas $\operatorname{AME}(4,d)$ has infinitely many LU classes for
  every $d\geq 4$ \sourcecite{ref:p43-rather}{RRKL23}.  The latter theorem
  includes infinitely many inequivalent $\operatorname{AME}(4,6)$ states, so
  neither side of this classification realizes
  Eq.~\eqref{eq:p43-finite-nontrivial-moduli}.

  \item Tan completely classifies five-qubit AME states: every such state is
  LU equivalent to a point of the unique $((5,2,3))$ code $\mathcal C$, and
  two points of $\mathcal C$ are equivalent exactly when related by its
  finite binary-tetrahedral transversal group
  \sourcecite{ref:p43-tan}{Tan26}.  Because the projective state space of the
  two-dimensional code is a continuum while this group is finite,
  $\lvert\mathfrak{M}_{5,2}\rvert$ is infinite, not finite and nontrivial.

  \item The 2026 AME review records the known one-class and infinite-class
  regimes and explicitly leaves Eq.~\eqref{eq:p43-finite-nontrivial-moduli}
  as open problem T5 \sourcecite{ref:p43-review}{RBR+26}.
\end{itemize}

\subsection*{References}

\begin{enumerate}[label={},leftmargin=4.8em,labelsep=0.5em]
  \footnotesize
  \item[\textup{[RRKL23]}]\label{ref:p43-rather}
  S. A. Rather, N. Ramadas, V. Kodiyalam, and A. Lakshminarayan,
  ``Absolutely Maximally Entangled State Equivalence and the Construction of
  Infinite Quantum Solutions to the Problem of 36 Officers of Euler,''
  \emph{Physical Review A} \textbf{108}, 032412 (2023).
  \href{https://doi.org/10.1103/PhysRevA.108.032412}{doi:10.1103/PhysRevA.108.032412};
  \href{https://arxiv.org/abs/2212.06737}{arXiv:2212.06737}.

  \item[\textup{[Tan26]}]\label{ref:p43-tan}
  I. Tan, ``Classification of Five-Qubit Absolutely Maximally Entangled
  States,'' arXiv:2507.02185v4 (2026).
  \href{https://arxiv.org/abs/2507.02185}{arXiv:2507.02185}.

  \item[\textup{[RBR+26]}]\label{ref:p43-review}
  G. Rajchel-Mieldzio{\'c}, R. Bistro{\'n}, A. Rico, A. Lakshminarayan, and
  K. {\.{Z}}yczkowski, ``Absolutely Maximally Entangled Pure States of
  Multipartite Quantum Systems,'' \emph{Reports on Progress in Physics}
  \textbf{89}, 057601 (2026).
  \href{https://doi.org/10.1088/1361-6633/ae672a}{doi:10.1088/1361-6633/ae672a};
  \href{https://arxiv.org/abs/2508.04777}{arXiv:2508.04777v3}.
\end{enumerate}

\subsection*{Comment}

The remaining task is either to exhibit an AME parameter pair whose full LU
quotient is a discrete non-singleton set, or to prove that every nonempty
$\mathfrak{M}_{N,d}$ is a singleton or infinite.  Classifying only a selected
construction family would not determine the full quotient in
Eq.~\eqref{eq:p43-finite-nontrivial-moduli}.

\subsection*{Field}

Quantum Resource Theory

\subsection*{Topic}

Absolutely maximally entangled states; Local unitary equivalence

\subsection*{ID}

\texttt{op\_0c86d9293aba3b01}
